\paragraph{Constant fibers supported on two roots.}
The following treats constant fibers supported on at most two roots in
every fixed value degree. It does not classify all repeated fibers.

\begin{lemma}[Three scaled two-root covers]
\label{lem:two-root-three-covers}
Let $k$ be algebraically closed of characteristic zero or characteristic
$p>r+s$, where $r,s\ge1$ are coprime. For nonconstant $q\in k(X)$,
at most two distinct $\lambda\in k^\times$ admit a rational solution of
\[
 Y^r(Y+1)^s=\lambda q(X).
\]
\end{lemma}
\begin{proof}
Put $b=r+s$, $h(Y)=Y^r(Y+1)^s$, and
$\kappa=(-1)^r r^rs^s/b^b$.
The cover $t=h(Y)/\lambda$ has inertia cycles $(r)(s)$ over zero,
a $b$-cycle over infinity, and a transposition over
$\kappa/\lambda$. Indeed
$h'=Y^{r-1}(Y+1)^{s-1}(bY+r)$; these are all branch points,
and ramification is tame. For $r=s=1$, zero is unramified.

Its monodromy group is $S_b$. To see primitivity, first exclude a
polynomial decomposition $h=A\circ B$ of degrees greater than one.
If $A$ has two distinct roots, their disjoint preimages under $B$
are supported on $\{0,-1\}$. Thus
$B-\alpha=cY^m$ and $B-\beta=c(Y+1)^m$, up to interchange,
whose difference cannot be constant for $m>1$ in the stated
characteristic. If $A$ has one root, its degree divides both $r$ and
$s$, also impossible. Any intermediate field gives such a polynomial
decomposition: by L\"uroth it is rational, and the unique point over
infinity makes both maps polynomial after choosing that point as
infinity. Hence the group is primitive. The graph whose edges are
conjugates of its transposition has connected components forming a
block system. It is therefore connected, and its edge transpositions
generate $S_b$.

For three distinct nonzero $\lambda_i$, take the compositum of the
three Galois closures. Its group is a subdirect subgroup of $S_b^3$.
Inertia at $\kappa/\lambda_i$ supplies a transposition in the $i$th
factor and identities in the others. Conjugation by the surjective
$i$th projection generates that entire factor, proving the group is
$S_b^3$. The normalized fiber product is thus connected of degree
$N=b^3$. At infinity it has $b^2$ inertia orbits; each of the three
unique simple branch points has index $b^2$. At zero the orbit count is
\[
 r^2+s^2+3r+3s.
\]
Here the all-$r$ and all-$s$ triples contribute $r^2,s^2$; the mixed
triples have orbit length $rs$, contributing $3r,3s$.
Riemann--Hurwitz gives
\[
 2g-2=2b^2-r^2-s^2-3b=b^2+2rs-3b
 \ge(b-2)(b+1)\ge0.
\]
Three rational solutions would embed this positive-genus function
field into $k(X)$ via $t=q(X)$, contradicting L\"uroth.
This also excludes inseparable $q$; no assumption $q'\ne0$ is used.
\end{proof}

\begin{theorem}[Two-root constant fibers]
\label{thm:two-root-constant-fiber}
Let $D\ge1$ and let $k$ be algebraically closed of characteristic zero or
characteristic $p>b=r+s$, with $r,s\ge1$. Suppose $H\in k[X]\setminus\{0\}$ and a polynomial
$F(X,u)$ has a constant fiber
\[
 F(X,u)-c_0H(X)=(u-R(X))^r(u-S(X))^s,
 \qquad \deg R,\deg S\le D.
\]
Put $d=\gcd(r,s)$. The degree-at-most-$D$ polynomial sections
$F(X,P)=cH$ either all lie in one affine polynomial pencil with
coefficient degrees at most $D$, or number at most $2b/d+2$.
Consequently, on any $n$ distinct evaluation points and for any received
word, at most
\[
 \max\left\{2b/d+2,\left\lfloor1/\eta\right\rfloor\right\}
\]
such sections agree in at least $D+\eta n$ coordinates, for $\eta>0$.
\end{theorem}
\begin{proof}
If $R\ne S$, put $a=R-S$, $Y=(P-R)/a$, $q=H/a^b$, and
$\lambda=c-c_0$. Then $Y^r(Y+1)^s=\lambda q$.
For nonconstant $q$ and $d=1$, Lemma~\ref{lem:two-root-three-covers}
allows at most two nonzero labels, each with at most $b$ rational
roots; the zero fiber adds only $Y=0,-1$.

For general $d$, write $h_0(Y)=Y^{r/d}(Y+1)^{s/d}$.
If a nonzero-label solution $Y_*$ exists, set $q_0=h_0(Y_*)$,
so $q=q_0^d/\lambda_*$. Every other nonzero-label solution has
\[
 \left(\frac{h_0(Y)}{q_0}\right)^d=\frac\lambda{\lambda_*},
\]
so $h_0(Y)=\mu q_0$ for a constant $\mu\ne0$.
The function $q_0$ is nonconstant. Applying the lemma to the coprime
exponents $r/d,s/d$ permits at most two $\mu$, each with at most
$b/d$ roots. Including the zero fiber proves $2b/d+2$.

If $q$ is constant, every $Y$ is constant, giving the pencil
$P=R+ta$. If $R=S$, the identity is
$(P-R)^b=(c-c_0)H$. One noncritical section $P_*$ expresses $H$
as a constant multiple of $(P_*-R)^b$, after which every section
belongs to the pencil $R+t(P_*-R)$. If none exists, there is only
$P=R$.

Finally, in a pencil $P_0+tP_1$ with $P_1\ne0$ of degree at most
$D$, at most $D$ coordinates have $P_1=0$; at every other coordinate
at most one member matches a fixed received symbol. If $L$ members
have at least $A>D$ agreements, then $LA\le LD+n$, proving
$L\le n/(A-D)$. Extending constants proves the same upper bound
for banks over smaller fields.
\end{proof}
